% =====================================================================
%  Kalender der FF KRZ -- Projektbeschreibung
%  Vorlage: deb-support (debsupport.sty, siehe latex/README.md)
%  Uebersetzen mit:  pdflatex projektbeschreibung_kalender_ff_krz.tex   (2x, wegen Seitenzahl)
% =====================================================================
\documentclass[a4paper,fontsize=10pt]{scrartcl}
\usepackage[bericht]{debsupport}

\usepackage{tikz}
\usetikzlibrary{positioning,arrows.meta,fit,backgrounds}

% Farben der Grafiken aus dem Corporate Design
\colorlet{ffrot}{dsrot}
\colorlet{grauhell}{dshell}
\colorlet{grau}{dsgrau}


\renewcommand*{\dsdokumentart}{Projektbeschreibung}
\renewcommand*{\dskurztitel}{Kalender der FF KRZ -- Projektbeschreibung}
\renewcommand*{\dsbeleg}{}
\setlength{\emergencystretch}{3em}

\begin{document}

\dstitel{Kalender der FF KRZ}
        {Automatische Bereitstellung der Termine aus MP Feuer auf Smartphones, Tablets und in Outlook}
        {Datum:   & \today\\
         Autor:   & \dsInhaber\\
         Projekt: & Kalender der FF KRZ\\
         Version: & Entwurf (Stand 12.11.2026)}

\tableofcontents

% ------------------------------------------------------------------
\section{Ausgangslage und Ziel}

Termine der Freiwilligen Feuerwehr KRZ werden in MP Feuer gepflegt, zum Beispiel
Übungsdienste, Sitzungen, Lehrgänge oder die Jahreshauptversammlung. Die
Mitglieder sehen diese Termine heute nicht automatisch in ihrem persönlichen
Kalender auf Smartphone, Tablet oder in Outlook.

Ziel des Projekts ist, \textbf{alle Termine aus MP Feuer automatisch und
ausschließlich lesend auf die Smartphones, Tablets und in das Outlook der
Mitglieder} zu bringen, und
zwar abhängig davon, welchen Einheiten, Gruppen und Abteilungen ein Mitglied
angehört.

Dazu gilt:
\begin{itemize}
  \item Für jede Einheit, Gruppe bzw.\ jeden Lehrgang gibt es einen eigenen Kalender, z.\,B.
        die Termine der LG Drove im Kalender \emph{LG Drove}, die Termine der
        JF Gr1 im Kalender \emph{JF Gr1} und die Termine des Lehrgangs TM1 2027
        im Kalender \emph{TM1 2027}.
  \item Jedes Mitglied bekommt genau die Kalender freigeschaltet, die zu
        seinen Zuordnungen und Mitgliedschaften passen.
  \item Die Kalender erscheinen in der normalen Kalender-App von Smartphone
        und Tablet (iPhone, iPad und Android) sowie in Outlook und werden
        automatisch aktuell gehalten.
  \item Termine werden nur in MP Feuer gepflegt. Auf den Endgeräten können
        die Kalender nur angesehen, nicht verändert werden.
\end{itemize}

% ------------------------------------------------------------------
\section{Systemüberblick und Informationsfluss}

Das System besteht aus drei Teilen:

\begin{enumerate}
  \item \textbf{MP Feuer} als führendes System. Hier werden Mitglieder,
        Zuordnungen und Termine wie bisher gepflegt.
  \item Ein \textbf{Kalenderserver}, der die exportierten Daten regelmäßig
        übernimmt, die Termine auf die Kalender verteilt und festlegt, welches
        Mitglied welche Kalender sehen darf. Er bietet ein einfaches
        Web-Frontend, auf dem jedes Mitglied seine Zugangsdaten, Kalender,
        Abo-Links und Mitgliedschaften einsehen kann, und verschickt
        Begrüßungs- und Hinweis-E-Mails.
  \item Die \textbf{Endgeräte der Mitglieder}: Smartphones und Tablets melden
        sich mit persönlichen Zugangsdaten am Kalenderserver an, Outlook
        abonniert die Kalender über persönliche Abo-Links. Beide rufen die
        Kalender automatisch ab.
\end{enumerate}

\begin{figure}[h]
\centering
\begin{tikzpicture}[
    font=\small,
    box/.style={draw, rounded corners=3pt, align=center, minimum height=1.1cm,
                text width=3.2cm, fill=white},
    data/.style={draw, align=center, text width=3.2cm, minimum height=0.9cm,
                 fill=grauhell},
    arr/.style={-{Stealth[length=2.5mm]}, thick},
    lbl/.style={font=\footnotesize, align=center}
]
  % Quelle
  \node[box, draw=ffrot, very thick] (mp) {\textbf{MP Feuer}\\Mitglieder, Gruppen,\\Termine};

  % Exporte
  \node[data, right=1.2cm of mp, yshift=0.9cm] (exm) {Export 1:\\Mitglieder und\\Mitgliedschaften};
  \node[data, right=1.2cm of mp, yshift=-0.9cm] (ext) {Export 2:\\Termine je\\Bereich/Gruppe};

  % Server
  \node[box, right=1.2cm of exm, yshift=-0.9cm, text width=3.6cm, minimum height=2.6cm] (srv)
       {\textbf{Kalenderserver}\\[3pt]
        Zuordnung\\Termin $\rightarrow$ Kalender\\[2pt]
        Freischaltung\\Mitglied $\rightarrow$ Kalender\\[2pt]
        Zugangsdaten\\[2pt]
        Web-Frontend, E-Mail};

  \draw[arr] (mp.east) -- ++(0.5,0) |- (exm.west);
  \draw[arr] (mp.east) -- ++(0.5,0) |- (ext.west);
  \draw[arr] (exm.east) -- ++(0.5,0) |- ([yshift=0.4cm]srv.west);
  \draw[arr] (ext.east) -- ++(0.5,0) |- ([yshift=-0.4cm]srv.west);

  \node[lbl, above=0.1cm of exm] {regelmäßig, automatisch};

  % Endgeräte
  \node[box, below=1.7cm of srv, xshift=-6.6cm, text width=3.2cm, minimum height=1.4cm] (ios) {\textbf{iPhone / iPad}\\eingebaute\\Kalender-App};
  \node[box, below=1.7cm of srv, xshift=-2.6cm, text width=3.2cm, minimum height=1.4cm] (and) {\textbf{Android}\\Smartphone/Tablet\\mit Sync-App};
  \node[box, below=1.7cm of srv, xshift=1.6cm, text width=3.2cm, minimum height=1.4cm] (out) {\textbf{Outlook}\\Abo-Link\\je Kalender};

  \draw[arr] (srv.south) -- ++(0,-0.7) -| (ios.north);
  \draw[arr] (srv.south) -- ++(0,-0.7) -| (and.north);
  \draw[arr] (srv.south) -- ++(0,-0.7) -| (out.north);
  \node[lbl, anchor=south west] at ([xshift=0.1cm,yshift=-0.65cm]srv.south) {nur lesend};

  \node[lbl, below=0.15cm of ios, xshift=2.0cm] {Benutzername und Passwort, verschlüsselt};
  \node[lbl, below=0.15cm of out] {geheimer persönlicher Link};
\end{tikzpicture}
\caption{Informationsfluss von MP Feuer bis zu den Endgeräten}
\end{figure}

Der Informationsfluss geht \textbf{immer nur in eine Richtung}: von MP Feuer
über den Kalenderserver zu den Endgeräten. Es fließen keine Daten zurück
nach MP Feuer.

Der Kalenderserver ist idealerweise unter einer \textbf{eigenen Domain}
erreichbar, z.\,B.\ \texttt{kalender.<domain-der-ff>}. Unter dieser Adresse
finden die Mitglieder das Web-Frontend, und dieselbe Adresse wird auf
Smartphone, Tablet und in Outlook eingetragen. Sie bleibt dauerhaft gleich,
auch wenn sich am Betrieb des Servers später etwas ändert. Die E-Mails des
Systems werden ebenfalls von dieser Domain verschickt. Die Domain muss vor dem
Echtbetrieb endgültig feststehen, weil sie in den Kennungen der Termine auf
den Geräten enthalten ist.

% ------------------------------------------------------------------
\section{Benötigte Exporte aus MP Feuer}

Damit der Kalenderserver arbeiten kann, braucht er zwei Exporte. Beide sollen
\textbf{regelmäßig und automatisch} erstellt werden, zum Beispiel mehrmals
täglich oder stündlich, damit Änderungen zeitnah ankommen.

\subsection*{Export 1: Mitglieder und Mitgliedschaften}

Eine Liste \textbf{aller Mitglieder der Feuerwehr}, mit:
\begin{itemize}
  \item eindeutiger Kennung des Mitglieds (z.\,B. Mitgliedsnummer). Diese
        Kennung darf sich nie ändern, auch nicht bei einer Namensänderung,
        da das Mitglied darüber wiedererkannt wird,
  \item Vor- und Nachname,
  \item E-Mail-Adresse (für Begrüßungs-E-Mail, Einmal-Links und
        Benachrichtigungen),
  \item allen Zuordnungen und Mitgliedschaften, also Einheit, Gruppen,
        Abteilungen und Lehrgänge (z.\,B.\ LG Drove, JF Gr1, KF, TM1 2027).
\end{itemize}

Berücksichtigt werden nur Mitglieder mit einer \textbf{offensichtlich gültigen
E-Mail-Adresse}: Die Adresse enthält ein „@“ und eine gültige Domain, also eine
Domain, die existiert und E-Mails empfangen kann. Diese Prüfung kann bereits
beim Export in MP Feuer erfolgen. Spätestens beim Import in den
Kalenderserver wird sie durchgeführt. Mitglieder ohne gültige Adresse erhalten
keinen Zugang. Sobald in MP Feuer eine gültige Adresse hinterlegt ist, wird
das Mitglied beim nächsten Import automatisch übernommen und erhält die
Begrüßungs-E-Mail.

Umgekehrt gilt: \textbf{Jedes Mitglied mit hinterlegter gültiger
E-Mail-Adresse erhält einen Zugang}, auch Kinder und Jugendliche. Mehrere
Mitglieder dürfen dieselbe Adresse haben, z.\,B.\ Geschwister in der
Kinderfeuerwehr mit der Adresse der Eltern. Sie erhalten getrennte Zugänge mit
eigenem Benutzernamen und Passwort.

Das gilt auch für Mitglieder, die bereits einen Zugang haben: Wird ihre Adresse
in MP Feuer ungültig, z.\,B.\ durch einen Tippfehler, wird der Zugang beim
nächsten Import gesperrt, bis die Adresse in MP Feuer korrigiert ist. So fallen
fehlerhafte Adressen schnell auf. Benutzername, Passwort und Abo-Links bleiben
dabei erhalten und funktionieren nach der Korrektur wieder.

\subsection*{Export 2: Termine}

Alle Termine der einzelnen Bereiche, Gruppen und Abteilungen, jeweils mit:
\begin{itemize}
  \item eindeutiger Kennung des Termins (damit Änderungen und Absagen
        erkannt werden),
  \item Bereich bzw.\ Gruppe, zu der der Termin gehört,
  \item Titel, Beginn und Ende (oder ganztägig),
  \item Ort und Beschreibung, soweit vorhanden,
  \item Status, z.\,B.\ ob ein Termin abgesagt wurde.
\end{itemize}

Ein Termin, der im Export nicht mehr enthalten ist, wird auch aus dem
Kalender entfernt.

% ------------------------------------------------------------------
\section{Verarbeitung im Kalenderserver}

Nach jedem Export führt der Kalenderserver folgende Schritte aus:

\begin{enumerate}
  \item \textbf{E-Mail-Adressen prüfen:} Mitglieder ohne offensichtlich
        gültige E-Mail-Adresse werden nicht übernommen (siehe Export~1).
  \item \textbf{Termine verteilen:} Jeder Termin wird anhand seines Bereichs
        dem passenden Kalender zugeordnet. Jeder Kalender existiert nur
        einmal, auch wenn viele Mitglieder ihn sehen.
  \item \textbf{Kalender freischalten:} Anhand der Mitgliedschaften wird für
        jedes Mitglied festgelegt, welche Kalender es sehen darf. Ändert sich
        eine Mitgliedschaft in MP Feuer, ändert sich auch die Freischaltung
        automatisch, ein Kalender kommt also hinzu oder fällt weg. Kalender,
        die alle Mitglieder sehen sollen (z.\,B.\ \emph{FF KRZ}), werden im
        Kalenderserver als \glqq für alle Mitglieder\grqq\ markiert und sind
        dann unabhängig von den Mitgliedschaften für jedes Mitglied
        freigeschaltet.
  \item \textbf{Zugänge verwalten:} Neue Mitglieder erhalten automatisch
        einen Zugang und eine Begrüßungs-E-Mail. Der
        Benutzername hat die Form \texttt{vorname.nachname}. Umlaute und ß
        werden dabei ausgeschrieben (z.\,B.\ \texttt{juergen.mueller}). Bei
        Namensgleichheit wird eine Ziffer angehängt (z.\,B.\
        \texttt{max.mustermann2}). Der Benutzername ändert sich danach nicht
        automatisch (siehe \glqq Namensänderung\grqq\ unten).
        Ändert sich die E-Mail-Adresse eines Mitglieds, wird die
        Begrüßungs-E-Mail an die neue Adresse geschickt. Passwort und
        Abo-Links bleiben dabei unverändert.
        Ausgeschiedene Mitglieder verlieren ihren Zugang automatisch; alle
        ihre Daten werden gelöscht.
  \item \textbf{Benachrichtigen:} Kommt für ein Mitglied ein Kalender hinzu,
        erhält es eine kurze E-Mail mit Hinweis auf den neuen Kalender und
        einem Anmeldelink für das Web-Frontend (14 Tage gültig). Den
        Abo-Link für Outlook findet es dort; Abo-Links werden nicht per
        E-Mail verschickt. Neue Mitglieder erhalten keinen zusätzlichen
        Hinweis, sondern nur die Begrüßungs-E-Mail mit ihrem Anmeldelink.
  \item \textbf{Betrieb protokollieren:} Zeitpunkt und Ergebnis jedes
        Imports werden festgehalten und im Web-Frontend angezeigt. Dazu
        gehört auch die Anzahl der Mitglieder, die wegen fehlender oder
        ungültiger E-Mail-Adresse nicht übernommen oder gesperrt wurden. Das
        vollständige Protokoll ist auf der Status-Seite für Administratoren
        einsehbar (siehe Abschnitt~\ref{sec:admin}).
\end{enumerate}

\textbf{Namensänderung:} Ändert sich der Name eines Mitglieds in MP Feuer,
z.\,B.\ durch Heirat, bleibt der Benutzername zunächst unverändert. Name und
E-Mail-Adresse werden im Kalenderserver aktualisiert, Passwort und Abo-Links
bleiben gleich. Das Mitglied muss an seinen Geräten nichts ändern.

\emph{Beispiel:} Anna Schmidt hat den Benutzernamen \texttt{anna.schmidt} und
nutzt ihn auf Smartphone und Tablet. Nach ihrer Heirat heißt sie in MP Feuer
Anna Becker. Ihr Benutzername bleibt \texttt{anna.schmidt}, und alle Geräte
funktionieren weiter. Der Kalenderserver erkennt sie nicht am Namen, sondern
an ihrer unveränderlichen Kennung aus MP Feuer. Ohne diese Kennung sähe es
so aus, als wäre Anna Schmidt ausgeschieden und Anna Becker neu
hinzugekommen. Ihr bisheriger Zugang würde gelöscht und ihre Geräte
funktionierten nicht mehr.

\textbf{Anzeige und Umstellung des Benutzernamens:} Der Kalenderserver
erkennt, wenn ein Benutzername nicht mehr zum aktuellen Namen passt. Das
Mitglied sieht dann im Web-Frontend einen Hinweis, die Administratoren sehen
es auf ihrer Status-Seite. Auf Wunsch kann der Zugang auf den passenden neuen
Benutzernamen umgestellt werden, entweder vom Mitglied selbst im
Web-Frontend oder von einem Administrator. Für die Umstellung gilt:

\begin{itemize}
  \item Sie ist nur möglich, wenn der Benutzername tatsächlich nicht mehr
        zum Namen passt, und muss ausdrücklich bestätigt werden.
  \item Der alte Benutzername ist danach ungültig und wird nicht an andere
        Mitglieder vergeben, solange das Mitglied in MP Feuer geführt wird.
        Scheidet es aus, werden alle seine Daten gelöscht; seine früheren
        Benutzernamen können dann wieder vergeben werden.
  \item Passwort, Abo-Links und Kalender bleiben unverändert. Outlook läuft
        ohne Anpassung weiter.
  \item Auf Smartphone und Tablet muss der Kalender-Account einmalig mit dem
        neuen Benutzernamen neu eingerichtet werden. Darauf wird vor der
        Bestätigung deutlich hingewiesen.
  \item Das Mitglied erhält eine E-Mail mit dem neuen Benutzernamen.
\end{itemize}

\emph{Im Beispiel:} Anna sieht im Web-Frontend den Hinweis, dass ihr
Benutzername \texttt{anna.schmidt} auf einem früheren Namen beruht. Sie kann
ihn auf \texttt{anna.becker} umstellen und richtet danach ihr Smartphone und
Tablet mit dem neuen Benutzernamen neu ein. Sie kann ihn aber auch
unverändert lassen.

Die Zuordnung von Bereichen in MP Feuer zu Kalendern wird einmalig
festgelegt und von den Administratoren im Web-Frontend gepflegt. Ein
Kalender kann dabei mehrere Bereiche umfassen. Beispiel:

\begin{center}
\begin{tabularx}{0.85\textwidth}{lX}
\toprule
\textbf{Kalender} & \textbf{Inhalt} \\
\midrule
FF KRZ   & Termine der gesamten Wehr, z.\,B.\ Jahreshauptversammlung
           (für alle Mitglieder) \\
LG Drove & Termine der Löschgruppe Drove \\
GJF      & Termine der Gemeindejugendfeuerwehr \\
JF Gr1   & Termine der Jugendfeuerwehr-Gruppe 1 \\
KF       & Termine der Kinderfeuerwehr \\
TM1 2027 & Termine des Lehrgangs TM1 2027 \\
\ldots   & weitere Einheiten, Gruppen und Abteilungen \\
\bottomrule
\end{tabularx}
\end{center}

\textbf{Beispiel:} Ein Mitglied der LG Drove, das zusätzlich in der
Jugendarbeit (GJF, JF Gr1) und in der Kinderfeuerwehr tätig ist, sieht auf
seinen Geräten die Kalender \emph{FF KRZ}, \emph{LG Drove}, \emph{GJF},
\emph{JF Gr1} und \emph{KF}. Jeder davon erscheint als eigener Kalender mit
eigener Farbe und kann einzeln ein- oder ausgeblendet werden.

% ------------------------------------------------------------------
\section{Bereitstellung auf den Endgeräten}

Es wird keine eigene App der Feuerwehr benötigt. Jedes Mitglied kann seinen
Zugang auf mehreren Geräten gleichzeitig einrichten, z.\,B.\ auf Smartphone,
Tablet und in Outlook. Die Inhalte sind überall gleich.

\subsection*{Smartphones und Tablets}

Smartphones und Tablets nutzen einen offenen, weit verbreiteten Standard für
die Kalendersynchronisation. Die Anmeldung erfolgt mit dem Benutzernamen
(\texttt{vorname.nachname}) und dem vom System festgelegten Passwort. Mit
diesem einen Zugang erhält das Gerät automatisch alle für das Mitglied
freigeschalteten Kalender.

\begin{itemize}
  \item \textbf{iPhone und iPad:} Die Kalender werden einmalig in den
        Einstellungen als zusätzlicher Kalender-Account eingerichtet, am
        einfachsten über ein Konfigurationsprofil aus dem Web-Frontend.
        Danach erscheinen sie in der eingebauten Kalender-App.
  \item \textbf{Android-Smartphones und -Tablets:} Android unterstützt den
        Standard nicht von Haus aus. Nötig ist einmalig eine verbreitete
        Synchronisations-App. Danach erscheinen die Kalender in jeder
        beliebigen Kalender-App.
  \item \textbf{Android ohne Zusatz-App:} Alternativ kann ein Mitglied die
        persönlichen Abo-Links (wie bei Outlook, siehe unten) im
        Google Kalender abonnieren. Die Kalender erscheinen dann in der
        Google-Kalender-App auf allen Geräten mit diesem Google-Konto.
        Google aktualisiert abonnierte Kalender allerdings nur alle paar
        Stunden; für kurzfristige Änderungen und Absagen ist die
        Synchronisations-App die bessere Wahl.
\end{itemize}

Die Geräte fragen den Kalenderserver in regelmäßigen Abständen nach
Änderungen, typischerweise alle 15 bis 60 Minuten, je nach Einstellung auf dem
Gerät. Eine Änderung in MP Feuer erscheint damit nach dem nächsten Export und
dem nächsten Abruf auf dem Gerät.

\subsection*{Outlook}

Outlook unterstützt den oben genannten Standard nicht von Haus aus. Es kann
aber in allen gängigen Varianten Internetkalender abonnieren: Outlook unter
Windows, das neue Outlook, Outlook im Web und Outlook auf dem Mac. Dafür stellt
der Kalenderserver zusätzlich \textbf{für jedes Mitglied und jeden
freigeschalteten Kalender einen persönlichen Abo-Link} bereit. Ein Abo-Link
ist nur so lange gültig, wie der Kalender dem Mitglied zugeordnet ist.

\begin{itemize}
  \item Das Mitglied abonniert in Outlook jeden seiner Kalender einzeln über
        den zugehörigen Link. Jeder Kalender erscheint dann als eigener
        Kalender und kann einzeln ein- oder ausgeblendet werden.
  \item Die Abo-Links stehen dem Mitglied ausschließlich im Web-Frontend zur
        Verfügung (siehe Abschnitt~\ref{sec:web}). Sie werden nicht per
        E-Mail verschickt.
  \item Der Kalender muss \textbf{abonniert}, nicht \textbf{importiert}
        werden. Ein Import (z.\,B.\ \glqq Aus Datei\grqq\ in Outlook,
        \glqq Importieren\grqq\ im Google Kalender oder das Öffnen einer
        heruntergeladenen \texttt{.ics}-Datei) kopiert die Termine nur
        einmalig; Änderungen und Absagen kämen nicht mehr an. Web-Frontend
        und Anleitungen weisen deutlich darauf hin. Optional zeigt der
        Kalenderserver eine Hinweisseite, wenn ein Abo-Link im Browser
        geöffnet wird, statt die Datei zum Herunterladen anzubieten.
        Smartphone und Tablet mit Benutzername und Passwort sind davon nicht
        betroffen.
  \item Outlook ruft abonnierte Kalender seltener ab als Smartphones,
        teilweise nur alle paar Stunden. Für kurzfristige Änderungen
        und Absagen sind Smartphone und Tablet daher die zuverlässigere
        Quelle.
  \item Kommt durch eine neue Mitgliedschaft ein Kalender hinzu, muss das
        Mitglied den neuen Link in Outlook einmalig zusätzlich abonnieren.
        Es erhält dazu eine E-Mail mit Anmeldelink und findet den neuen
        Abo-Link im Web-Frontend.
        Fällt die Zuordnung zu einem Kalender weg, wird der zugehörige Link
        automatisch ungültig und liefert keine Termine mehr.
\end{itemize}

Dieselben Abo-Links funktionieren auch im \textbf{Google Kalender}
(Abonnieren \glqq per URL\grqq\ auf calendar.google.com), mit denselben
Einschränkungen: Jeder Kalender wird einzeln abonniert, und Google ruft die
Kalender nur selten ab.

\subsection*{Einrichtung}

Für die Einrichtung wird je eine kurze Anleitung für iPhone/iPad, Android
(mit Synchronisations-App oder über den Google Kalender) und Outlook
bereitgestellt, zunächst als Text, nach Inbetriebnahme des Servers
bebildert.

% ------------------------------------------------------------------
\section{Web-Frontend und E-Mail-Versand}\label{sec:web}

Der Kalenderserver bietet ein einfaches Web-Frontend. Das Web-Frontend dient
nur der Information. Termine und Mitgliedschaften werden
weiterhin ausschließlich in MP Feuer gepflegt.

\subsection*{Anmeldung}

Für die Anmeldung am Web-Frontend gibt es zwei Wege:

\begin{itemize}
  \item \textbf{Mit Passwort:} Benutzername (\texttt{vorname.nachname})
        \emph{oder} E-Mail-Adresse, jeweils zusammen mit dem vom System
        festgelegten Passwort. Es ist dasselbe Passwort wie auf Smartphone
        und Tablet. Teilen sich mehrere Mitglieder eine E-Mail-Adresse,
        bestimmt das Passwort, welcher Zugang angemeldet wird.
  \item \textbf{Mit Einmal-Link:} Das Mitglied klickt auf einen persönlichen
        Anmeldelink aus einer E-Mail des Systems, z.\,B.\ aus der
        Begrüßungs-E-Mail. Ein neuer Link kann jederzeit auf der Anmeldeseite
        angefordert werden; bei mehreren Zugängen mit derselben Adresse
        erhält jeder Zugang einen eigenen Link. Jeder Link ist nur einmal
        verwendbar und zeitlich begrenzt gültig.
\end{itemize}

Der Einmal-Link ist auch der Weg, wenn das Passwort vergessen wurde: Nach der
Anmeldung per Link wird das Passwort im Web-Frontend angezeigt.

\subsection*{Inhalte des Web-Frontends}

\begin{itemize}
  \item \textbf{Meine Kalender} (Startseite): schlichte Liste aller für das
        Mitglied freigeschalteten Kalender mit Farbe, Beschreibung und der
        Anzahl der Termine im laufenden Kalenderjahr. Bei jedem Kalender
        steht, aus welcher Mitgliedschaft in MP Feuer er stammt oder dass er
        für alle Mitglieder gilt; Mitgliedschaften ohne eigenen Kalender
        werden darunter genannt. So lässt sich nachvollziehen, warum ein
        Kalender angezeigt wird oder fehlt. Technische Angaben wie Abo-Links
        oder Passwörter stehen bewusst nicht hier.
  \item \textbf{Kalender abonnieren:} eigener Menüpunkt mit zwei
        Unterseiten:
        \begin{itemize}
          \item \emph{Smartphone \& Tablet} (iPhone, iPad, Android):
                Serveradresse, Benutzername und Passwort sowie das Profil für
                iPhone/iPad. Das Passwort wird \emph{ausschließlich} hier
                angezeigt.
          \item \emph{Outlook \& Google Kalender:} zu jedem Kalender der
                persönliche Abo-Link zum Kopieren, mit kurzer Erklärung.
        \end{itemize}
        Dazu die Anleitungen für iPhone/iPad, Android (auch ohne Zusatz-App
        über den Google Kalender) und Outlook.
  \item \textbf{Terminübersicht als PDF:} Ein Jahreskalender zum Ausdrucken
        oder Aufhängen mit den eigenen Kalendern, wahlweise auf einer oder
        zwei Seiten (siehe unten).
  \item \textbf{Benutzername:} Hinweis, wenn der Benutzername nicht mehr
        zum aktuellen Namen passt, mit der Möglichkeit, ihn umzustellen
        (siehe \glqq Namensänderung\grqq\ in Abschnitt~4).
  \item \textbf{Neu generieren:} Das Mitglied kann jederzeit selbst ein
        neues Passwort und neue Abo-Links erzeugen lassen, z.\,B.\ wenn ein
        Gerät verloren wurde oder ein Link weitergegeben wurde. Das bisherige
        Passwort bzw.\ die bisherigen Links werden dabei sofort ungültig und
        müssen auf den Geräten bzw.\ in Outlook ersetzt werden. Beide
        Aktionen stehen auf der jeweiligen Unterseite von \glqq Kalender
        abonnieren\grqq.
  \item \textbf{Letzte Abrufe:} Zu jedem Abo-Link und zum Zugang für
        Smartphone und Tablet wird angezeigt, wann zuletzt erfolgreich
        Termine abgerufen wurden. So sieht das Mitglied selbst, ob ein Gerät
        oder Outlook noch aktualisiert, z.\,B.\ nach einem Passwortwechsel.
  \item \textbf{Betriebsinformationen:} Zeitpunkt des letzten erfolgreichen
        Imports aus MP Feuer und ob dieser fehlerfrei war, damit erkennbar
        ist, wie aktuell die Daten sind.
\end{itemize}

\subsection*{Terminübersicht als PDF}

Jedes Mitglied kann sich im Web-Frontend eine Terminübersicht als PDF
herunterladen, z.\,B.\ für den Kühlschrank oder die Pinnwand im Gerätehaus.

\begin{itemize}
  \item \textbf{Auswahl:} Startmonat (Vorgabe: Januar des laufenden Jahres),
        Zeitraum von 1 bis 12 Monaten (Vorgabe: 12) und die gewünschten
        Kalender. Angeboten werden nur die eigenen Kalender.
  \item \textbf{Aufteilung:} alle Monate auf einer Seite A4 quer oder auf zwei
        Seiten (z.\,B.\ Januar--Juni und Juli--Dezember) mit größeren Feldern.
  \item \textbf{Darstellung:} je Monat eine Spalte, je Tag eine Zeile. Jeder
        Termin erscheint als farbiges Oval in der Farbe seines Kalenders mit
        der Startzeit darin, ganztägige Termine als leeres Oval. Passen nicht
        alle Termine eines Tages, zeigt das Feld nur farbige Punkte. Abgesagte
        Termine erscheinen als leerer Kreis mit farbigem Rand.
  \item \textbf{Markierungen:} Wochenenden und gesetzliche Feiertage in NRW
        sind grau hinterlegt, Schulferien in NRW als Streifen am Rand.
  \item \textbf{Bewusst knapp:} Titel, Orte und Beschreibungen stehen nicht im
        PDF. Die Einzelheiten sieht das Mitglied wie gewohnt auf Smartphone,
        Tablet oder in Outlook. Ein ausgedrucktes oder weitergegebenes PDF
        verrät damit kaum etwas.
  \item Das PDF ist eine Momentaufnahme. Der Zeitpunkt der Erstellung steht
        in der Fußzeile.
\end{itemize}

Administratoren haben keine Ansicht der Termine; die Terminübersicht steht nur
den Mitgliedern für ihre eigenen Kalender zur Verfügung.

\subsection*{E-Mail-Versand}

Der Kalenderserver verschickt automatisch E-Mails an die Mitglieder.
Passwörter und Abo-Links werden grundsätzlich \emph{nicht} per E-Mail
verschickt.

\begin{itemize}
  \item \textbf{Begrüßungs-E-Mail} an jedes neue Mitglied sowie an jede
        neue oder geänderte E-Mail-Adresse. Sie erklärt kurz das System und
        enthält den Benutzernamen, einen Einmal-Link zur direkten Anmeldung am
        Web-Frontend (14 Tage gültig) sowie Hinweise auf die Anleitungen. Ein
        Passwort ist \emph{nicht} enthalten, es wird nur im Web-Frontend
        angezeigt.
  \item \textbf{Einmal-Link} für das Web-Frontend auf Anforderung über die
        Anmeldeseite (30 Minuten gültig).
  \item \textbf{Hinweis auf neue Kalender}, wenn durch eine neue
        Mitgliedschaft ein Kalender hinzukommt, mit einem Anmeldelink für das
        Web-Frontend (14 Tage gültig). Neue Mitglieder erhalten stattdessen
        nur die Begrüßungs-E-Mail.
  \item \textbf{Neuer Benutzername}, nachdem der Benutzername umgestellt
        wurde, mit Hinweis auf die Neueinrichtung von Smartphone und Tablet.
\end{itemize}

Die E-Mail-Adressen stammen aus dem Export von MP Feuer. Der Kalenderserver
erkennt bei jedem Import neue Mitglieder und geänderte E-Mail-Adressen und
verschickt die Begrüßungs-E-Mail dann automatisch. Bei einer geänderten Adresse
bleiben Benutzername, Passwort und Abo-Links unverändert. Bereits
eingerichtete Geräte und Outlook funktionieren also ohne Anpassung weiter.

% ------------------------------------------------------------------
\section{Status-Seite für Administratoren}\label{sec:admin}

Für die Betreuer des Systems gibt es eine eigene Status-Seite im
Web-Frontend. Sie ist nur für Benutzer mit Administratorrechten sichtbar und
dient der Überwachung des Betriebs.

\begin{itemize}
  \item \textbf{Letzter Import:} Zeitpunkt, Dauer und Ergebnis
        (erfolgreich oder fehlerhaft) des letzten Imports aus MP Feuer sowie
        eine Übersicht der vorherigen Importe.
  \item \textbf{Importprotokoll:} Was sich beim Import geändert hat, unter
        anderem:
        \begin{itemize}
          \item neu angelegte Mitglieder,
          \item ausgeschiedene Mitglieder,
          \item Mitglieder, die wegen fehlender oder ungültiger E-Mail-Adresse
                nicht übernommen oder gesperrt wurden, mit Namen, damit die
                Adresse in MP Feuer gezielt korrigiert werden kann,
          \item geänderte E-Mail-Adressen und Mitgliedschaften,
          \item Anzahl neuer, geänderter und entfernter Termine je Kalender,
          \item Fehler und Warnungen, z.\,B.\ Termine, die keinem Kalender
                zugeordnet werden konnten.
        \end{itemize}
  \item \textbf{Warnungen per E-Mail:} Ist ein Import fehlgeschlagen oder
        liegt seit mehr als 24 Stunden kein Export vor, erhalten alle
        Administratoren sowie zusätzlich hinterlegte Adressen eine E-Mail.
  \item \textbf{Benutzernamen:} Liste der Mitglieder, deren Benutzername
        nicht mehr zum aktuellen Namen passt, jeweils mit altem
        Benutzernamen, aktuellem Namen und der Möglichkeit, den Zugang auf
        den neuen Benutzernamen umzustellen.
  \item \textbf{Gesperrte Anmeldungen:} Nach zu vielen fehlgeschlagenen
        Anmeldeversuchen wird die Anmeldung für einen Benutzer bzw.\ eine
        IP-Adresse für bis zu eine Stunde gesperrt. Die Status-Seite listet
        alle aktuellen Sperren mit Anzahl der Versuche und Ende der Sperre.
        Ein Administrator kann jede Sperre einzeln oder alle auf einmal
        aufheben, z.\,B.\ wenn am Dienstabend im Gerätehaus viele Geräte
        gleichzeitig eingerichtet werden. IP-Adressen werden nur gekürzt
        angezeigt.
  \item \textbf{E-Mail-Versand:} Welche E-Mails verschickt wurden und ob der
        Versand fehlgeschlagen ist.
  \item \textbf{Nutzung insgesamt:} Wie viele Mitglieder ihren Zugang auf
        Smartphone, Tablet oder in Outlook aktiv nutzen, also in einem
        bestimmten Zeitraum erfolgreich Termine abgerufen haben.
  \item \textbf{Kalender und Zuordnungen:} Pflege der Kalender (Name, Farbe,
        \glqq für alle Mitglieder\grqq) und der Zuordnung von Bereichen aus
        MP Feuer zu Kalendern.
\end{itemize}

\subsection*{Nutzerübersicht}

Eine Tabelle über alle Mitglieder zeigt, wer das System bereits nutzt. Je
Mitglied:

\begin{itemize}
  \item Name, Benutzername und Status (aktiv oder wegen ungültiger
        E-Mail-Adresse gesperrt),
  \item wann die Begrüßungs-E-Mail verschickt wurde und ob bzw.\ wann der
        Link darin aufgerufen wurde,
  \item erste und letzte Anmeldung im Web-Frontend,
  \item letzter erfolgreicher Abruf über Smartphone bzw.\ Tablet, mit
        Gerätetyp,
  \item letzter erfolgreicher Abruf über einen Outlook-Abo-Link,
  \item die freigeschalteten Kalender mit Namen und Farbe. Bei vielen
        Kalendern zeigt die Liste die ersten drei und \glqq +n\grqq\ für den
        Rest. Termininhalte werden nicht angezeigt.
\end{itemize}

Die Übersicht lässt sich sortieren und filtern, z.\,B.\ nach
\glqq hat sich noch nie angemeldet\grqq, \glqq Begrüßungslink nie
aufgerufen\grqq, \glqq kein Abruf seit mehr als 30 Tagen\grqq\ oder
\glqq gesperrt\grqq. Zusätzlich lässt sie sich nach einem Kalender filtern,
z.\,B.\ \glqq alle Mitglieder mit Kalender LG Drove\grqq. So sehen die
Administratoren, wen sie gezielt ansprechen sollten, und können prüfen, ob die
Zuordnung von Bereichen zu Kalendern stimmt.

\subsection*{Detailansicht eines Mitglieds}

Ein Klick auf den Namen in der Nutzerübersicht öffnet die Detailansicht des
Mitglieds. Sie zeigt auf einen Blick:

\begin{itemize}
  \item Name, Benutzername, Status und die in MP Feuer hinterlegte
        E-Mail-Adresse,
  \item alle freigeschalteten Kalender mit dem letzten Abruf je Abo-Link,
  \item die Mitgliedschaften laut MP Feuer, aus denen sich die Kalender
        ergeben,
  \item die Nutzung: Begrüßung verschickt, Begrüßungslink aufgerufen,
        erste und letzte Anmeldung, letzter Abruf über Smartphone bzw.\
        Tablet und über Outlook,
  \item eine aktuelle Anmeldesperre wegen zu vieler Fehlversuche, mit der
        Möglichkeit, sie aufzuheben,
  \item die Aktionen \glqq Begrüßungs-E-Mail erneut senden\grqq\ und, wenn
        der Benutzername nicht mehr zum Namen passt, \glqq Benutzername
        umstellen\grqq.
\end{itemize}

Die Detailansicht hilft vor allem bei Rückfragen, etwa \glqq Ich sehe den
Kalender der JF nicht\grqq: Ist der Kalender nicht freigeschaltet, liegt es an
der Mitgliedschaft in MP Feuer und nicht am Gerät. Über \glqq Zurück zur
Nutzerübersicht\grqq\ geht es zur Liste zurück.

\subsection*{Begrüßungs-E-Mail erneut senden}

Hat ein Mitglied die Begrüßungs-E-Mail nicht beachtet oder gelöscht, kann ein
Administrator sie in der Detailansicht des Mitglieds mit einem Klick erneut
senden. Sie enthält dann einen neuen, wieder 14 Tage gültigen Einmal-Link;
ältere, nicht genutzte Begrüßungslinks werden dabei ungültig. Passwort und
Abo-Links bleiben unverändert.

Unabhängig davon kann sich jedes Mitglied jederzeit selbst über die
Anmeldeseite einen neuen Einmal-Link an seine E-Mail-Adresse schicken lassen.
Dafür muss es nur die Adresse des Kalenderservers kennen. Diese sollte daher
zusätzlich dauerhaft bekannt gemacht werden, z.\,B.\ auf der Website der
Feuerwehr oder per Aushang im Gerätehaus.

% ------------------------------------------------------------------
\section{Sicherheit und Datenschutz}

\begin{itemize}
  \item \textbf{Nur lesend:} Auf den Endgeräten können Termine nicht
        angelegt, geändert oder gelöscht werden. MP Feuer bleibt die einzige
        Stelle, an der Termine gepflegt werden.
  \item \textbf{Persönlicher Zugang:} Jedes Mitglied meldet sich mit eigenem
        Benutzernamen und Passwort an und sieht ausschließlich die für es
        freigeschalteten Kalender. Passwörter werden vom Kalenderserver
        zufällig erzeugt und nicht vom Mitglied selbst gewählt. Sie werden
        nicht per E-Mail verschickt, sondern ausschließlich nach Anmeldung im
        Web-Frontend angezeigt. Das Mitglied kann dort jederzeit ein neues
        Passwort erzeugen lassen. Das alte Passwort wird dann sofort
        ungültig. Passwörter aus MP Feuer werden dafür \emph{nicht}
        verwendet oder exportiert.
  \item \textbf{Begrenzung von Anmeldeversuchen:} Nach zu vielen
        Fehlversuchen wird die Anmeldung je Benutzer bzw.\ IP-Adresse für
        bis zu eine Stunde gesperrt. Administratoren sehen die Sperren auf der
        Status-Seite und können sie aufheben.
  \item \textbf{Einmal-Links:} Einmal-Links für das Web-Frontend werden nur
        an die in MP Feuer hinterlegte E-Mail-Adresse geschickt, sind
        zeitlich begrenzt gültig und nur einmal verwendbar. Sie werden erst
        durch einen Klick auf \glqq Jetzt anmelden\grqq\ verbraucht, damit
        Link-Prüfer von Mailprogrammen sie nicht vorab entwerten. Sie
        berechtigen nur zur Anmeldung am Web-Frontend, nicht zum Abruf der
        Kalender auf den Geräten.
  \item \textbf{Abo-Links für Outlook:} Da Outlook bei Abos die Anmeldung mit
        Benutzername und Passwort nicht zuverlässig unterstützt, ist jeder Abo-Link lang,
        zufällig und persönlich. Er gilt nur für ein Mitglied und einen
        Kalender und nur, solange dem Mitglied dieser Kalender zugeordnet
        ist. Bei Austritt oder Wegfall der Zuordnung wird er automatisch
        ungültig. Abo-Links werden nur im Web-Frontend angezeigt, nicht per
        E-Mail verschickt. Das Mitglied kann seine Abo-Links im Web-Frontend
        jederzeit neu erzeugen lassen. Die bisherigen Links werden dann sofort
        ungültig.
  \item \textbf{Verschlüsselte Übertragung:} Die Verbindung zwischen
        Endgerät und Kalenderserver ist durchgehend verschlüsselt.
  \item \textbf{Datensparsamkeit:} Aus MP Feuer werden nur die Daten
        übernommen, die für die Kalender nötig sind. Weitere
        Personaldaten werden nicht exportiert.
  \item \textbf{Nutzungsdaten:} Gespeichert werden nur Zeitpunkte (letzte
        Anmeldung, letzter Abruf, Aufruf des Begrüßungslinks) und der
        Gerätetyp, kein Verlauf einzelner Zugriffe. Nutzerübersicht und
        Detailansicht sind nur für Administratoren sichtbar. Sie zeigen die
        Namen der freigeschalteten Kalender, aber keine Termininhalte. Die Mitglieder werden in der
        Begrüßungs-E-Mail und im Web-Frontend darauf hingewiesen.
  \item \textbf{Automatisches Entziehen von Rechten:} Scheidet ein Mitglied
        aus oder wechselt es die Gruppe, werden Zugang bzw.\ Kalender beim
        nächsten Export automatisch angepasst. Ausgeschiedene Mitglieder
        werden mit allen Daten gelöscht, einschließlich früherer
        Benutzernamen.
\end{itemize}

% ------------------------------------------------------------------
\section{Mögliche Erweiterungen}

\subsection*{Kalenderansicht im Web-Frontend}

Zusätzlich zur Terminübersicht als PDF kann das Web-Frontend später eine
Kalenderansicht für die Mitglieder erhalten: eine Monatsansicht und eine
Liste der nächsten Termine, jeweils mit Titel, Zeit und Ort, nur für die
eigenen Kalender und ebenfalls nur lesend.

\subsection*{Kontakte}

Mit demselben Prinzip können später auch \textbf{Kontakte} verteilt werden.
Jedes Mitglied erhält dann abhängig von seinen Gruppen ein Adressbuch, z.\,B.\
mit den Ansprechpartnern seiner Einheit, ebenfalls nur lesend und automatisch
aktuell. Dafür wäre ein dritter Export mit den freizugebenden Kontaktdaten
nötig.

% ------------------------------------------------------------------
\section{Offene Punkte und nächste Schritte}

\begin{enumerate}
  \item Klären, ob und in welcher Form MP Feuer die beiden Exporte
        automatisch und regelmäßig bereitstellen kann (Format, Häufigkeit,
        Übertragungsweg, Zeitraum der Termine).
  \item Klären, ob Termine und Mitgliedschaften in MP Feuer dieselben
        Schlüssel für Bereiche verwenden.
  \item Festlegen, welche Bereiche, Gruppen und Abteilungen in MP Feuer zu
        welchem Kalender gehören, und die Liste der Kalender abstimmen.
  \item Fehlende oder ungültige E-Mail-Adressen in MP Feuer ergänzen, damit
        möglichst alle Mitglieder Zugang erhalten.
  \item Abstimmung mit der bzw.\ dem Datenschutzbeauftragten und der
        Wehrleitung.
  \item Eigene Domain für den Kalenderserver festlegen und einrichten,
        einschließlich Absenderadresse für die E-Mails des Systems.
  \item Klären, wo der Kalenderserver betrieben wird und wer ihn betreut.
  \item Testbetrieb mit einer kleinen Gruppe, danach schrittweise Einführung
        für alle Mitglieder.
\end{enumerate}

\end{document}
